\chapter{A proof of Euclidean coarea from area}\label{sec:euclidean-coarea-from-area}
We develop the regular-map and critical-map steps in \cite{Sim18}, Chapter 2, Theorem 7.3, printed pp.~76--79, and then pass to Lipschitz maps by the $C^1$ approximation theorem. The Euclidean area formula is now proved in Theorem \ref{thm:euclidean-area-local-proof}; Whitney extension remains an external input. The preliminary slicing estimate is proved from diameter covers; in particular it does not use the coarea theorem being established. All Hausdorff measures use the normalization in Definition \ref{def:hausdorff-measure}.

\begin{lemma}[Measurability of the measure of a compact fiber]\label{lem:compact-fiber-measure-borel}
Let $K\subset\mathbb R^N$ be compact, let $f:K\to\mathbb R^k$ be continuous, and let $d$ be an integer with $1\leq d\leq N$. Then
\[
 y\longmapsto\mathcal H^d(K\cap f^{-1}(y))
\]
is a Borel function with values in $[0,\infty]$.
\end{lemma}
\begin{proof}
For each positive integer $j$, define $c_j(y)$ as the infimum of $\sum_i\omega_d(\operatorname{diam}O_i/2)^d$ over finite covers of the compact fiber by open sets $O_i$ of diameter less than $1/j$. Restrict to the countable family of finite unions of balls with rational centers and radii. There are countably many finite cover candidates. A candidate covers the fiber precisely when $y\notin f(K\setminus\bigcup_iO_i)$. The set on the right being excluded is compact, so the set of levels where the candidate is a cover is open. Consequently $\{y:c_j(y)<a\}$ is open for every $a$, and $c_j$ is Borel.

The family restrictions do not change the limit: $c_j(y)\geq\mathcal H^d_{1/j}(K\cap f^{-1}(y))$. Conversely, start with a diameter cover of the fiber by sets of diameter less than $1/(2j)$, with total cost within $\varepsilon$ of the corresponding Hausdorff content. Enlarge each set slightly to an open finite union of rational balls, increasing its cost by at most $\varepsilon2^{-i}$ and keeping its diameter below $1/j$. Such an enlargement exists by a finite fine grid cover of each bounded set. Compactness of the fiber selects a finite subcover from these open sets. Thus $c_j(y)\leq\mathcal H^d_{1/(2j)}(K\cap f^{-1}(y))+2\varepsilon$ whenever the latter content is finite. Infinite content causes no upper-bound restriction. Sending $\varepsilon\downarrow0$ and then $j\to\infty$ proves that the increasing functions $c_j$ converge to the Hausdorff measure of the fiber. Their pointwise limit is Borel. Empty fibers have cost zero throughout.
\end{proof}

\begin{lemma}[Measurable fibers for a finite Hausdorff-measure set]\label{lem:finite-hausdorff-fiber-measurable}
Let $1\leq k<m\leq N$, let $A\subset U\subset\mathbb R^N$ be $\mathcal H^m$-measurable with finite measure, with $U$ open, and let $f:U\to\mathbb R^k$ be locally Lipschitz. Then $y\mapsto\mathcal H^{m-k}(A\cap f^{-1}(y))$, read as outer measure at exceptional levels if needed, is Lebesgue measurable. Almost every fiber is Hausdorff measurable.
\end{lemma}
\begin{proof}
Choose a Borel core of $A$ by Lemma \ref{lem:hausdorff-borel-core}. Its finite restricted measure is Radon, so inner regularity gives increasing compact sets $K_j\subset A$ with $\mathcal H^m(A\setminus\bigcup_jK_j)=0$. Lemma \ref{lem:compact-fiber-measure-borel} shows that each compact fiber measure is Borel. Continuity from below of Hausdorff measure on the increasing compact fibers gives a measurable fiber measure for $\bigcup_jK_j$.

The omitted set has null fibers for almost every level by Lemma \ref{lem:null-set-lipschitz-fibers}. This lemma was proved directly by covers, without coarea. At the remaining levels the full fiber differs from a countable union of compact sets by a set of zero Hausdorff outer measure. It is therefore Hausdorff measurable and has the same measure as that union. Since Lebesgue measure is complete, modifying the measurable fiber function on the exceptional null set preserves measurability. The same conclusion holds for sigma-finite sets by disjoint finite-measure pieces and a countable union of exceptional level sets.
\end{proof}

\begin{proposition}[A coarse Lipschitz slicing estimate]\label{prop:coarse-lipschitz-slicing}
Let $1\leq k<m\leq N$, let $f:\mathbb R^N\to\mathbb R^k$ have Lipschitz constant $L$, and let $A$ be $\mathcal H^m$-measurable with finite measure. Then
\[
 \int_{\mathbb R^k}\mathcal H^{m-k}(A\cap f^{-1}(y))\,dy
 \leq \frac{2^m\omega_k\omega_{m-k}}{\omega_m}L^k\mathcal H^m(A).
\]
The constant is not sharp. This is a diameter-cover version of \cite{Sim18}, Chapter 2, Theorem 1.10, pp.~49--51; it is sufficient for the critical-set argument below.
\end{proposition}
\begin{proof}
Put $d=m-k>0$. Choose diameter covers $E_{li}$ of $A$ with positive diameters $s_{li}<1/l$ and normalized $m$-cost at most $\mathcal H^m(A)+1/l$. Countably many zero-diameter members have zero $d$-measure in every fiber and can be discarded. For each remaining member meeting $A$, select $x_{li}\in A\cap E_{li}$ and replace it by the closed ball $Q_{li}=\overline B_{s_{li}}(x_{li})$. These balls still cover $A$ apart from the harmless countable set. Their images under $f$ are compact and lie in target balls of radius $Ls_{li}$. The measurable function
\[
 F_l(y)=\sum_i\omega_d s_{li}^d\chi_{f(Q_{li})}(y)
\]
therefore satisfies
\[
 \int F_l\,dy\leq\omega_d\omega_kL^k\sum_i s_{li}^{m}
 \leq\frac{2^m\omega_d\omega_k}{\omega_m}L^k
       \bigl(\mathcal H^m(A)+1/l\bigr).
\]
The balls contributing to $F_l(y)$ cover the fiber and have diameter less than $2/l$, so its Hausdorff content at that scale is at most $F_l(y)$. Hence the fiber measure is bounded by $\liminf_lF_l(y)$. It is measurable by Lemma \ref{lem:finite-hausdorff-fiber-measurable}. Fatou's lemma gives the stated integral bound. The factor $2^m$ records the passage from diameter sets to enclosing balls; no sharp isodiametric inequality is being asserted. For $L=0$, every image is a point of zero target Lebesgue measure and the same argument gives zero.
\end{proof}

\begin{proposition}[Coarea for a C1 map of full rank]\label{prop:full-rank-c1-coarea}
Let $U\subset\mathbb R^n$ be open and let $f:U\to\mathbb R^k$ be $C^1$, with $1\leq k<n$ and $\operatorname{rank}Df=k$ everywhere. For every Lebesgue-measurable $A\subset U$,
\[
 \int_AJ_kDf\,dx=\int_{\mathbb R^k}\mathcal H^{n-k}(A\cap f^{-1}(y))\,dy.
\]
This regular case uses only ordinary change of variables, Fubini, Euclidean area on the fibers and the linear Jacobian identity.
\end{proposition}
\begin{proof}
Near each point, choose $k$ input coordinates on which the corresponding minor of $Df$ is invertible. The map $g(x)=(f(x),p(x))$, with $p$ the remaining $n-k$ coordinates, is a $C^1$ diffeomorphism on a sufficiently small neighborhood. Write its inverse as $\psi(y,z)$ on an open set $W\subset\mathbb R^k\times\mathbb R^{n-k}$. On a fixed level, $z\mapsto\psi(y,z)$ is injective and regular. The identity $f\circ\psi(y,z)=y$ and Lemma \ref{lem:coarea-chart-jacobian-identity} give
\[
 J_kDf(\psi(y,z))\,|\det D\psi(y,z)|
 =J_{n-k}(D_z\psi(y,z)).
\]
Indeed use $A=D\psi$, $B=(I_k,0)$ and $\ker B=\{0\}\times\mathbb R^{n-k}$ in that lemma.

Ordinary change of variables on the diffeomorphism and Fubini now give, for a Borel subset $A$ of this neighborhood,
\[
 \int_AJ_kDf\,dx
 =\int_{\mathbb R^k}\int_{\{z:(y,z)\in W\}}
 \chi_A(\psi(y,z))J_{n-k}(D_z\psi)\,dz\,dy.
\]
The injective case of Theorem \ref{thm:area-formula} identifies the inner integral with the fiber measure. Cover $U$ by countably many such neighborhoods and partition a Borel $A$ among them. Additivity on the disjoint pieces and monotone convergence prove the formula for Borel sets, allowing infinite values. For general Lebesgue-measurable $A$, choose a Borel representative modulo a null set on a countable bounded exhaustion. Lemma \ref{lem:null-set-lipschitz-fibers} removes the remainder on almost every fiber. This proves the general assertion without using a coarea theorem for the regular case.
\end{proof}

\begin{lemma}[Adding small target directions makes a critical map regular]\label{lem:critical-coarea-regularization-jacobian}
Let $U\subset\mathbb R^n$ be open and let $f:U\to\mathbb R^k$ be $C^1$, with $1\leq k<n$. For $0<\varepsilon\leq1$, define
\[
 F_\varepsilon(x,z)=f(x)+\varepsilon z,
 \qquad (x,z)\in U\times(0,1)^k.
\]
This map has rank $k$ everywhere. If $\operatorname{rank}Df(x)<k$ and $\|Df(x)\|_{\mathrm{op}}\leq L$, then
\[
 J_kDF_\varepsilon(x,z)
 =\sqrt{\det(Df(x)Df(x)^*+\varepsilon^2I_k)}
 \leq\varepsilon(L^2+1)^{(k-1)/2}.
\]
This is the regularization in \cite{Sim18}, proof of Theorem 7.3, pp.~78--79.
\end{lemma}
\begin{proof}
The derivative matrix is $[Df,\varepsilon I_k]$, so its row rank is $k$ and its Gram matrix is the one displayed. At a critical point the nonnegative eigenvalues of $DfDf^*$ lie in $[0,L^2]$ and at least one is zero. After adding $\varepsilon^2I_k$, at least one eigenvalue is exactly $\varepsilon^2$ and the other $k-1$ eigenvalues are at most $L^2+1$. Their product and its square root give the bound. For $k=1$ the estimate is exactly $J_1DF_\varepsilon=\varepsilon$ on the critical set.
\end{proof}

\begin{proposition}[The critical set of a C1 map has null almost-everywhere fibers]\label{prop:c1-critical-fibers-null}
Let $f:U\subset\mathbb R^n\to\mathbb R^k$ be $C^1$, with $U$ open and $1\leq k<n$. Put $C=\{x:\operatorname{rank}Df(x)<k\}$. Then
\[
 \mathcal H^{n-k}(C\cap f^{-1}(y))=0
 \quad\text{for almost every }y\in\mathbb R^k.
\]
This conclusion is proved without assuming the Euclidean coarea formula for critical maps.
\end{proposition}
\begin{proof}
It suffices to prove the assertion for a Borel subset $A\subset C$ in a closed ball whose concentric larger closed ball lies in $U$; countably many such subsets cover $C$. The derivative has a finite bound $L$ on the larger ball. Integration along segments in that ball shows that $f$ is Lipschitz on the smaller ball, so the fiber-measurability lemma applies to $A$, and $|A|<\infty$. Apply Proposition \ref{prop:full-rank-c1-coarea} to $F_\varepsilon$ from Lemma \ref{lem:critical-coarea-regularization-jacobian}, on the measurable set $A\times(0,1)^k$. Writing
\[
 B_y=(A\times(0,1)^k)\cap F_\varepsilon^{-1}(y),
\]
we obtain
\[
 \int_{\mathbb R^k}\mathcal H^n(B_y)\,dy
 \leq\varepsilon(L^2+1)^{(k-1)/2}|A|.
\]
In particular $B_y$ is Borel and has finite $n$-measure at almost every level. Apply Proposition \ref{prop:coarse-lipschitz-slicing} on each such set to the one-Lipschitz projection $(x,z)\mapsto z$. Its fiber at $z\in(0,1)^k$ is an isometric copy of $A\cap f^{-1}(y-\varepsilon z)$. Thus, with $C_{n,k}=2^n\omega_k\omega_{n-k}/\omega_n$ and $H(v)=\mathcal H^{n-k}(A\cap f^{-1}(v))$,
\[
 \int_{(0,1)^k}H(y-\varepsilon z)\,dz
 \leq C_{n,k}\mathcal H^n(B_y)
 \quad\text{for almost every }y.
\]
The function $H$ is Lebesgue measurable by Lemma \ref{lem:finite-hausdorff-fiber-measurable}. A Borel representative may be used in the double integral: contain its exceptional set in a Borel null set, whose translate is null on each fixed $z$ slice, and Tonelli shows that the joint exceptional set has product measure zero. Integrate the last inequality in $y$ and change variables $v=y-\varepsilon z$. Since the cube has volume one, this gives
\[
 \int_{\mathbb R^k}H(v)\,dv
 \leq C_{n,k}\varepsilon(L^2+1)^{(k-1)/2}|A|.
\]
The left side is independent of $\varepsilon$. Let it tend to zero; nonnegativity forces $H=0$ almost everywhere. Countable compact localization completes the proof. The coarse slicing estimate was proved for arbitrary finite-measure sets, so its use on $B_y$ presupposes neither rectifiability of $A$ nor a critical coarea theorem.
\end{proof}

\begin{theorem}[Euclidean coarea for a C1 map]\label{thm:c1-euclidean-coarea-local}
For a $C^1$ map $f:U\to\mathbb R^k$ on an open $U\subset\mathbb R^n$, $1\leq k<n$, and a Lebesgue-measurable $A\subset U$,
\[
 \int_AJ_kDf\,dx=\int_{\mathbb R^k}\mathcal H^{n-k}(A\cap f^{-1}(y))\,dy.
\]
Both sides can be infinite. The proof retains Euclidean area and ordinary Euclidean integration as inputs.
\end{theorem}
\begin{proof}
The regular set $U\setminus C$ is open, and Proposition \ref{prop:full-rank-c1-coarea} proves the formula on $A\setminus C$. On $A\cap C$, the Jacobian is zero. Proposition \ref{prop:c1-critical-fibers-null} gives zero fiber measure at almost every level, so both sides for this part vanish. The two pieces are disjoint; adding them proves the formula. Measurability on a countable bounded exhaustion is Lemma \ref{lem:finite-hausdorff-fiber-measurable}. The decomposition does not claim that the critical values form a null set; only the relevant lower-dimensional measure of almost every critical fiber is zero.
\end{proof}

\begin{proposition}[C1 agreement reduces a Lipschitz map to smooth coarea]\label{prop:lipschitz-coarea-c1-reduction}
Let $f:\mathbb R^n\to\mathbb R^k$ be globally Lipschitz, with $1\leq k<n$, and let $A\subset\mathbb R^n$ be bounded and Lebesgue measurable. Then
\[
 \int_AJ_kDf\,dx=\int_{\mathbb R^k}\mathcal H^{n-k}(A\cap f^{-1}(y))\,dy.
\]
Here $Df$ is supplied by Rademacher and the Jacobian is zero on its exceptional set. The proof uses Whitney-based $C^1$ approximation and the preceding $C^1$ formula.
\end{proposition}
\begin{proof}
For every positive integer $j$, Corollary \ref{cor:vector-c1-lusin-approximation} gives a $C^1$ map $g_j$ agreeing with both $f$ and its derivative off a set of measure less than $1/j$. Let $E_j$ be the differentiability points where both equalities hold. Their union has full measure, since its complement lies in each exceptional set and thus has measure at most $1/j$ for every $j$. Make the $E_j$ disjoint by assigning each point to the first qualifying index. These agreement sets are measurable because $f,g_j$ and the derivative representatives are measurable.

Apply Theorem \ref{thm:c1-euclidean-coarea-local} to $g_j$ on $A\cap E_j$. On this set $J_kDg_j=J_kDf$, and its fiber under $g_j$ is exactly its fiber under $f$. Sum the identities over $j$. Additivity of measure on the disjoint fiber pieces and monotone convergence identify the sum with the coarea identity for $A\cap\bigcup_jE_j$. The omitted subset has zero volume, and Lemma \ref{lem:null-set-lipschitz-fibers} gives zero lower-dimensional fiber measure at almost every level. It affects neither side. This proves the formula without using Euclidean coarea in the $C^1$ approximation or null-fiber steps.
\end{proof}
